\documentclass[a4paper]{article}
% Shared preamble for the English probability and statistics documents.
% Based on the original course preamble. Compile with pdfLaTeX.
\usepackage[top=1.5cm,bottom=2cm,left=2cm,right=2cm]{geometry}
\usepackage[T1]{fontenc}
\usepackage{lmodern}
\usepackage[utf8]{inputenc}
\usepackage[english]{babel}

\usepackage{amsmath}
\usepackage{amssymb}
\usepackage{amsthm}
\usepackage{mathrsfs}
\usepackage{mathtools}
\usepackage{etoolbox}
\usepackage{tikz}
\usetikzlibrary{shapes.multipart,arrows.meta}
\usepackage{xcolor}
\usepackage{enumerate}
\usepackage[shortlabels]{enumitem}
\usepackage{hyperref}
\usepackage{version}
\usepackage{bbold}
\usepackage{fontawesome5}
\usepackage{fancyhdr}
\lfoot{}
\usepackage{natbib}
\usepackage{graphicx}
\usepackage{pdfpages}
\usepackage{authblk}
    
\usepackage{blindtext}
\usepackage{multicol}
\usepackage{collcell}
\usepackage{wrapfig}
\usepackage{blkarray}
\usepackage{cancel}
\usepackage{appendix}
\usepackage{float}
\makeatletter
\input{tex/assets/slashbox.sty}
\makeatother
\hypersetup{unicode,colorlinks,citecolor=red,filecolor=black,
  linkcolor=red,urlcolor=cyan}
\setlength{\columnsep}{1cm}
\pagenumbering{gobble}

\def\numtext#1{%
  \ifcase#1\or one\or two\or three\or four\or five\or six\or seven\or eight\or nine\or ten\or
  eleven\or twelve\or thirteen\or fourteen\or fifteen\or sixteen\or seventeen\or eighteen\or nineteen\or twenty\or Lots\fi}

\newcommand{\Byte}[1]{%
  \def\nodecontents{}%
  \foreach[count=\l] \k in {#1} {\xappto{\nodecontents}{\noexpand\nodepart{\numtext{\l}}\k}}%
  \node[rectangle split,rectangle split parts=8,rectangle split horizontal,draw] at (2,2){\nodecontents};}

\newcommand{\D}{{\mathrm{d}}}
\newcommand{\cE}{{\mathcal{E}}}
\newcommand{\cG}{{\mathcal{G}}}
\newcommand{\cX}{{\mathcal{X}}}
\newcommand{\cY}{{\mathcal{Y}}}
\newcommand{\cB}{{\mathcal{B}}}
\newcommand{\cF}{{\mathcal{F}}}
\newcommand{\cO}{{\mathcal{O}}}
\newcommand{\cS}{{\mathcal{S}}}
\newcommand{\cT}{{\mathcal{T}}}
\newcommand{\given}{\, \big| \,}
\newcommand{\cL}{{\mathscr{L}}}
\newcommand{\cD}{{\mathcal{D}}}
\newcommand{\cA}{{\mathscr{A}}}
\newcommand{\cP}{{\mathscr{P}}}
\newcommand{\cR}{{\mathscr{R}}}
\newcommand{\A}{{\mathcal A}}
\newcommand{\one}{\mathbb{1}}
\newcommand{\mathbbm}[1]{\mathbb{#1}}

\providecommand{\C}{}
\renewcommand{\C}{\mathbb C}
\newcommand{\R}{\mathbb R}
\newcommand{\Q}{\mathbb Q}
\newcommand{\N}{\mathbb N}
\newcommand{\T}{\mathbb T}
\newcommand{\F}{\mathbb F}
\renewcommand{\P}{\mathbb P}
\newcommand{\Z}{\mathbb Z}
\newcommand{\E}{\mathbb {E}}
\newcommand{\Pbb}[1]{\P\left( #1\right)}
\newcommand{\Ebb}[1]{\E\left[ #1\right]}

\newcommand{\lb}{\left (}
\newcommand{\rb}{\right )}
\newcommand{\lbb}{\left [}
\newcommand{\rbb}{\right ]}
\newcommand{\labs}{\left |}
\newcommand{\rabs}{\right |}
\newcommand{\lbrb}[1]{\lb #1 \rb}
\newcommand{\lbbrbb}[1]{\lbb#1\rbb}
\newcommand{\labsrabs}[1]{\labs#1\rabs}
\newcommand{\lbbrb}[1]{\lbb#1\rb}
\newcommand{\lbrbb}[1]{\lb#1\rbb}
\newcommand{\langlerangle}[1]{\langle#1\rangle}
\newcommand{\lbcurly}{\{}
\newcommand{\rbcurly}{\}}
\newcommand{\curly}[1]{\lbcurly#1\rbcurly}
\newcommand{\abs}{\labsrabs}
\newcommand{\ind}[1]{1_{#1}}
\newcommand{\indep}{\perp \!\!\! \perp}
\renewcommand{\d}{\, \mathrm{d}}

\newtheorem{theorem}{Theorem}
\newtheorem{problemCnt}{Lemma}
\theoremstyle{definition}
\newtheorem{problem}[problemCnt]{Problem}
\newtheorem{problembook}[problemCnt]{\faBook \,  Problem}
\newtheorem{problempencil}[problemCnt]{\faPencil* \,  Problem}
\newtheorem{problemstar}[problemCnt]{\faStar\,  Problem}
\newtheorem{problemstarpencil}[problemCnt]{\faStar \, \faPencil* \,   Problem}
\newtheorem{probleminteresting}[problemCnt]{\faSmile[regular] \,   Problem}
\newtheorem{solution}{Solution}

\newcommand{\fwcell}[1]{\makebox[\arraystretch\normalbaselineskip]{$#1$}}
\newenvironment{enumeratei}{\begin{enumerate}[(i)]}{\end{enumerate}}
\newenvironment{enumeratea}{\begin{enumerate}[(a)]}{\end{enumerate}}
\newenvironment{enumeraten}{\begin{enumerate}[\upshape 1.]}{\end{enumerate}}

% Compatibility helpers needed by a few files in the combined archive.
\AtBeginDocument{\noindent{\footnotesize English translation prepared with ChatGPT (OpenAI, GPT-6), 26 October 2026, for translation of the Bulgarian text into English and LaTeX formatting.
The Bulgarian originals remain the reference; translation may contain errors.}\par\smallskip}
\newcommand{\MissingGraphic}[1]{%
  \begin{center}\fbox{\parbox{0.82\linewidth}{\centering Missing source image: \texttt{#1}}}\end{center}}



\setlength{\emergencystretch}{3em}
\pagenumbering{arabic}
\begin{document}
\begin{center}\Large Probability and Statistics, Computer Science\par\large Test 1, 13 May 2023: solutions\end{center}
\noindent\textit{English transcription of the supplied handwritten workings. The notation and derivations follow the notes, including inconsistencies identified below.}
\section*{Problem 1}
\begin{enumerate}
\item Number the prize boxes from 0 to 4. The reached box number is $Y\sim Bin(4,1/2)$. For the payout $X$,
\[\begin{array}{c|ccccc}X&1000&100&700&250&500\\\hline\P&1/16&4/16&6/16&4/16&1/16\end{array}\]
\[\E X=\frac{1000+400+4200+1000+500}{16}=\frac{7100}{16}=443.75.\]
\item Let $X_1$ be the number of balls until the first 1000-lev prize. Then $X_1\sim Ge(1/16)$ and $\E X_1=16$. Let $X_2$ be the total payout until then. Informally, the final 1000 leva are preceded by an average of 15 other payouts. Their conditional law is
\[\begin{array}{c|cccc}X\mid X\ne1000&100&700&250&500\\\hline\P&4/15&6/15&4/15&1/15\end{array}\]
so $\E[X\mid X\ne1000]=6100/15=406\frac23$, yielding $1000+6100=7100$.
Formally,
\begin{align*}
\E X_2&=\sum_{k\ge1}\P(X_1=k)\E[X_2\mid X_1=k]\\
&=\sum_{k\ge1}\left(\frac{15}{16}\right)^{k-1}\frac1{16}\left[(k-1)\frac{6100}{15}+1000\right].
\end{align*}
Use
\[\sum_{k=1}^\infty\left(\frac{15}{16}\right)^{k-1}(k-1)=\sum_{j=0}^\infty j\left(\frac{15}{16}\right)^j=\frac{15/16}{(1/16)^2}=15\cdot16\]
and $\sum_{k\ge1}(15/16)^{k-1}/16=1$. Hence $\E X_2=7100$.
\item Let $Z_i$ be the randomly assigned prize in box $i$. Each is uniform on $\{1000,100,700,250,500\}$, so $\E Z_i=2550/5=510$. Then
\[\E[\text{payout}]=\E\sum_{i=0}^4Z_i\binom4i(1/2)^4=510.\]
The proposed best arrangement in the notes is $100,700,1000,500,250$.
\item The probability of hitting the 100-lev prize on each of three drops is
\[f(p)=(4p(1-p)^3)^3=64p^3(1-p)^9.\]
\[f'(p)=192p^2(1-p)^9-576p^3(1-p)^8.\]
The interior critical point satisfies $192(1-p)=576p$, giving $p=1/4$.
\item Bayes' rule gives
\[\frac{(3/8)^2}{(3/8)^2+(8/27)^2}\approx62\%.\]
\end{enumerate}
\section*{Problem 2}
Let $Y_1=X_1+X_2$ and $Y_2=X_1X_2$, with parameters indexed by 1 or 2 as appropriate.
\subsection*{1. Bernoulli variables}
For iid $X_i\sim Ber(p)$, the sum takes values in $\{0,1,2\}$ and is not Bernoulli. Since $\P(X_1X_2=1)=p^2$, the product is $Ber(p^2)$.
\subsection*{2. Binomial variables}
For iid $X_i\sim Bin(n,p)$, the sum of the two independent sums of $n$ Bernoulli variables is $Bin(2n,p)$.
If $X_1X_2\sim Bin(n_2,p_2)$, its largest possible value forces $n_2=n^2$. Its expectation gives
\[\E X_1X_2=n^2p^2=n_2p_2,\qquad p_2=p^2.\]
The probability of its maximum would have to satisfy
\[p^{2n}=\P(X_1X_2=n^2)=p_2^{n_2}=p^{2n^2}.\]
The notes conclude a contradiction except for the degenerate $p=1$ case, where $X_1X_2\sim Bin(n^2,1)$. (The $n=1$ Bernoulli case is treated separately above.)
\subsection*{3. Geometric variables}
The notes use the geometric convention with mean $1/p$. If $X_1+X_2\sim Ge(p_2)$, its mean would give $2/p=1/p_2$, hence $p_2=p/2$. Matching variances would require
\[\frac{2(1-p)}{p^2}=\frac{1-p_2}{p_2^2}=\frac{4-2p}{p^2},\]
which is impossible.
If $X_1X_2\sim Ge(p_3)$, the mean gives $p_3=p^2$. On the other hand,
\begin{align*}
D(X_1X_2)&=\E X_1^2X_2^2-(\E X_1X_2)^2\\
&=\left(\frac{1-p}{p^2}+\frac1{p^2}\right)^2-\left(\frac1{p^2}\right)^2
=\frac{3-4p+p^2}{p^4}.
\end{align*}
Equating this with $(1-p^2)/p^4$ gives $2p^2-4p+2=0$, or $p=1$. Thus only the degenerate $Ge(1)$ case works.
\subsection*{4. Poisson variables}
For iid $X_i\sim Po(\lambda)$, $X_1+X_2\sim Po(2\lambda)$, for example by generating functions (see the exercises). If the product were $Po(\lambda_2)$, its mean would force $\lambda_2=\lambda^2$, but
\[D(X_1X_2)=(\lambda+\lambda^2)^2-\lambda^4=\lambda^2+2\lambda^3.\]
Matching mean and variance would give $2\lambda^3=0$, a contradiction for $\lambda>0$.
\section*{Problem 3}
Write
\[f_\sigma=\tfrac14\bigl(f(X_1,X_2,X_3)+f(X_2,X_1,X_3)+f(X_3,X_2,X_1)+f(X_1,X_3,X_2)\bigr).\]
\subsection*{1. Sum}
For $f=X_1+X_2+X_3$, $\E f=3\mu$, $Df=3\sigma^2$, and $f_\sigma=f$. Thus both means and variances agree.
\subsection*{2. Signed sum}
For $f=X_1+X_2-X_3$,
\[f_\sigma=\tfrac14\bigl(X_1+X_2-X_3+X_3+X_2-X_1+X_2+X_1-X_3+X_1+X_3-X_2\bigr)=\frac{X_1+X_2}{2}.\]
Hence $\E f_\sigma=\E f=\mu$, while $Df=3\sigma^2$ and
\[Df_\sigma=\tfrac14(DX_1+DX_2)=\sigma^2/2.\]
The averaging has an effect: the expectation is unchanged and the variance is smaller.
\subsection*{3. Product minus the third variable}
The notes define
\[Y_1=X_1X_2-X_3,\quad Y_2=X_3X_2-X_1,\quad Y_3=X_2X_3-X_1=Y_2,\quad Y_4=X_1X_3-X_2.\]
For each, $\E Y_i=\mu^2-\mu$, and
\begin{align*}
Df=DY_1&=D(X_1X_2)+DX_3\\
&=(\sigma^2+\mu^2)^2-\mu^4+\sigma^2
=\sigma^2(\sigma^2+2\mu^2+1).
\end{align*}
The subsequent handwritten calculation is
\begin{align*}
Df_\sigma&=\tfrac1{16}D(Y_1+2Y_2+Y_4)\\
&=\tfrac38DY_1+\tfrac58Cov(Y_1,Y_2),\\
Df-Df_\sigma&=\tfrac58\bigl(DY_1-Cov(Y_1,Y_2)\bigr).
\end{align*}
If the latter is positive, $f_\sigma$ has smaller variance. Expanding the covariance gives
\begin{align*}
Cov(X_1X_2-X_3,X_3X_2-X_1)
&=Cov(X_1X_2,X_2X_3)-2Cov(X_1,X_1X_2)\\
&=\mu^2(\sigma^2+\mu^2)-\mu^4-2\bigl((\sigma^2+\mu^2)\mu-\mu^3\bigr)\\
&=\sigma^2(\mu^2-2\mu).
\end{align*}
Therefore the displayed conclusion is
\[Df-Df_\sigma=\tfrac58\sigma^2(\sigma^2+\mu^2+2\mu+1)>0.\]
\noindent\textit{Transcription note: this part of the handwriting duplicates $Y_2$ as $Y_3$, although the preceding definition of $f_\sigma$ uses a different permutation. It also writes $\E f_\sigma=\mu$ on the following page, inconsistent with its earlier $\mu^2-\mu$. These source inconsistencies have not been silently replaced by a new solution.}
\end{document}
